\section{Installing RAFTcorr}
\label{sec:installing}

Pick the route that fits your computer. Every route gives you the same application, the same
\file{raftcorr} command line and the same model files.

\begin{center}
\begin{tabularx}{\linewidth}{@{}p{0.4\linewidth}X@{}}
  \toprule
  You have & Use \\
  \midrule
  Windows 10 (version 1809 or later) or 11, 64-bit & the Setup program, \file{RAFTcorr-Windows-Setup.exe}
    (Section~\ref{sec:installing-windows}) \\
  A Mac with Apple silicon, macOS 14 or later & the disk image, \file{RAFTcorr-macOS.dmg}
    (Section~\ref{sec:installing-macos}) \\
  Python 3.10 or later, on any system & \file{pip install raftcorr}
    (Section~\ref{sec:installing-pip}) \\
  A copy of the source code & \file{install.bat} on Windows, or a conda environment
    (Section~\ref{sec:installing-source}) \\
  \bottomrule
\end{tabularx}
\end{center}

Both downloads are on the release page, \url{https://github.com/zachtong/RAFTcorr/releases/latest}.

\subsection{Windows: the Setup program}
\label{sec:installing-windows}

The Setup program is small: it holds a private Python and RAFTcorr itself, but not PyTorch, the
library that runs the models. Which PyTorch build fits depends on your graphics card, so Setup
chooses and downloads it while it installs. It needs an internet connection once, for PyTorch,
the other Python packages and the model files.

\begin{enumerate}
  \item Download \file{RAFTcorr-Windows-Setup.exe} from the release page and run it. It installs
    for your own user account and asks for no administrator password, so it works on a lab
    computer too. Windows may warn you first; see the note below.
  \item Accept the licence and keep the proposed folder,
    \file{\%LOCALAPPDATA\%\textbackslash Programs\textbackslash RAFTcorr}, or choose another one.
  \item Choose the additional tasks. \ui{Add the "raftcorr" command to PATH (new terminals)} is
    ticked by default; \ui{Create a desktop shortcut} is not. Click \ui{Install}.
  \item Setup copies its files, then opens a console window and installs the Python components in
    three steps: \file{[1/3] PyTorch}, \file{[2/3] RAFTcorr} and its dependencies,
    \file{[3/3] model files}. The status line reads \emph{``Installing PyTorch, RAFTcorr and the
    model files (up to 3 GB with a GPU; this takes a while)...''}. Wait for the console window to
    close by itself.
  \item On the last page, keep \ui{Launch RAFTcorr} ticked to start the application. The box is
    offered only when the Python components were installed.
\end{enumerate}

\begin{warnbox}
Do not close the console window while it works: that stops the installation of the Python
components. If it happened, run \ui{Repair RAFTcorr} from the Start menu.
\end{warnbox}

\subsubsection*{Which PyTorch}
Setup asks the NVIDIA driver (the program \file{nvidia-smi}) which CUDA version it supports and
installs the matching PyTorch build:

\begin{center}
\begin{tabular}{ll}
  \toprule
  The driver reports & PyTorch build \\
  \midrule
  CUDA 12.8 or later, or a GeForce RTX 50 series card & CUDA 12.8 \\
  CUDA 12.0 to 12.7 & CUDA 12.6 \\
  CUDA 11.x & CUDA 11.8 \\
  CUDA 10 or older, no CUDA version, or no NVIDIA driver & CPU build \\
  \bottomrule
\end{tabular}
\end{center}

The console line of step \file{[1/3]} says which case applied, for example
\emph{``[1/3] PyTorch: driver supports CUDA 12.4: PyTorch with CUDA 12.6''}. Without an NVIDIA GPU, RAFTcorr runs
on the CPU, more slowly. If the driver is too old, the line ends with \emph{``(update the NVIDIA
driver for the GPU)''}: update the driver, then run \ui{Repair RAFTcorr}.

\subsubsection*{What you get}
\begin{itemize}
  \item \ui{RAFTcorr} in the Start menu, and on the desktop if you asked for it.
  \item \ui{Repair RAFTcorr} in the Start menu. It runs the second half of the setup again:
    detects the driver, installs PyTorch, RAFTcorr and the model files. A PyTorch of the build
    that fits the driver is kept, so nothing is downloaded for it again; a PyTorch of another
    build is uninstalled first and the right one installed, for example the CUDA build in place
    of the CPU build after a driver update. Use it when the setup was interrupted (no connection,
    console window closed) or after you installed or updated an NVIDIA driver.
  \item The \file{raftcorr} command in every new terminal, when the PATH task was ticked. Only the
    folder \file{bin} of the installation goes on the PATH; the private Python and its
    \file{pip} stay off it.
  \item \file{.raftcorr} project files open in RAFTcorr when you double-click them.
  \item \file{install.log} in the installation folder: the full output of every setup and
    repair, the PyTorch build chosen and, at the end, the versions found
    (\file{raftcorr info}, Section~\ref{sec:installing-check}).
\end{itemize}

If the Python components could not be installed, Setup says so: \emph{``RAFTcorr could not install
its Python components (code~N). Check the internet connection, then run `Repair RAFTcorr' from the
Start menu.''}, followed by the path of \file{install.log}. A model file that could not be
downloaded does not stop the setup; RAFTcorr fetches the model when it is first used.

\begin{notebox}
RAFTcorr is not code-signed, so Windows SmartScreen may stop the Setup program the first time
with \emph{``Windows protected your PC''}. Click \ui{More info}, then \ui{Run anyway}. SmartScreen
reports that the file is not yet widely downloaded; it is not a virus detection.
\end{notebox}

\subsubsection*{Upgrading}
Run the Setup program of the new version; there is no need to uninstall first. It installs into
the same folder, replaces the private Python with a clean one and installs PyTorch again, so it
needs the internet connection again. Your model files, settings and projects stay. Close RAFTcorr
before you start.

\subsubsection*{Uninstalling}
Open the Windows settings, \ui{Apps}, then \ui{Installed apps} (Windows~11) or
\ui{Apps \& features} (Windows~10), and uninstall \ui{RAFTcorr}. This removes the installation
folder with its private Python and \file{install.log}, the PATH entry and the \file{.raftcorr}
file association. The model files in
\file{\%LOCALAPPDATA\%\textbackslash RAFTcorr\textbackslash models} and your settings stay;
delete the folder by hand if you no longer need the models.

\subsection{macOS: the disk image}
\label{sec:installing-macos}

\begin{enumerate}
  \item Download \file{RAFTcorr-macOS.dmg} from the release page.
  \item Open it and drag \ui{RAFTcorr} onto \ui{Applications}.
  \item Start RAFTcorr from Applications, Launchpad or Spotlight.
\end{enumerate}

It needs a Mac with Apple silicon and macOS 14 (Sonoma) or later; Intel Macs are not supported.
The application is self-contained: Python, PyTorch and Qt are inside, and nothing else is
installed. PyTorch runs on the CPU, since a Mac has no CUDA. The model files are not inside; they
are downloaded the first time each is used (Section~\ref{sec:installing-models}). To upgrade,
drag the new version onto \ui{Applications} and replace the old one.

\begin{notebox}
The application is not notarised by Apple, so the first time macOS refuses to open it. Open
\ui{System Settings}, then \ui{Privacy \& Security}, scroll to the line about RAFTcorr, click
\ui{Open Anyway} and confirm. From then on it opens like any other application. A new version may
ask again.
\end{notebox}

The command line is inside the application. To use it as \file{raftcorr}, add an alias to your
shell, for example in \file{\textasciitilde/.zshrc}:

\begin{lstlisting}
alias raftcorr=/Applications/RAFTcorr.app/Contents/MacOS/RAFTcorr
\end{lstlisting}

Without arguments this command opens the application; with a command such as \file{run} or
\file{models} it is the command line of Section~\ref{sec:cli}. RAFTcorr does not register
\file{.raftcorr} project files on macOS: open a project with \menu{File}{Open Project…}.

\subsection{With pip}
\label{sec:installing-pip}

Use this route on Linux, or when you work in Python already. It needs Python 3.10 or later.
Install RAFTcorr into a virtual environment of its own, so its dependencies cannot collide with
other packages.

On a computer with an NVIDIA GPU, install PyTorch from the PyTorch index first, choosing the CUDA
build that your driver supports (the table of Section~\ref{sec:installing-windows};
\url{https://pytorch.org/get-started/locally/} lists the builds). The PyTorch on PyPI is the CPU
build on Windows, and \file{pip install raftcorr} alone would install it. On a Mac, skip this
step.

\begin{lstlisting}
pip install torch --index-url https://download.pytorch.org/whl/cu126
pip install raftcorr
raftcorr                     # start the application
raftcorr models --download   # fetch the model files now (otherwise on first use)
raftcorr selftest            # run every installed model on a synthetic pair
\end{lstlisting}

\file{python -m raftcorr} works like \file{raftcorr} when the script folder of pip is not on your
PATH. pip also installs \file{raftcorr-gui}, which starts the application without a console
window on Windows.

\begin{notebox}
On Windows, PySide6 and PyTorch install files up to about 150 characters deep below
\file{site-packages}. If pip stops with \emph{``No such file or directory''} and a hint about long
paths, create the environment in a shorter folder, for example
\file{C:\textbackslash venvs\textbackslash raftcorr}, or enable long paths in Windows
(\url{https://pip.pypa.io/warnings/enable-long-paths}).
\end{notebox}

If PyTorch is already installed as the CPU build, \file{pip install torch -{}-index-url ...} keeps
it. Remove it first with \file{pip uninstall torch}, then install the CUDA build.

\subsection{From the source code}
\label{sec:installing-source}

Clone the repository, \url{https://github.com/zachtong/RAFTcorr}. On Windows with Anaconda or
Miniconda, run \file{install.bat} in its folder. It creates the conda environment
\file{raftcorr-app} (Python 3.11), detects the GPU as the Setup program does, installs the
matching PyTorch build (a download of about 2.5\,GB), RAFTcorr in editable mode and the model
files, and writes \file{install.log} in the same folder. Like Repair, it keeps a PyTorch of the
matching build, so running it again downloads no PyTorch, and uninstalls a PyTorch of another
build before it installs the matching one. Start the application with
\file{run.bat}.

On any platform:

\begin{lstlisting}
conda env create -f environment.yml
conda activate raftcorr-app
pip install torch --index-url https://download.pytorch.org/whl/cu126  # or /cpu
pip install -e ".[dev]"
\end{lstlisting}

In a source checkout, RAFTcorr also looks for model files in its folder \file{models}.

\subsection{First launch}
\label{sec:installing-first}

Start RAFTcorr from the Start menu, Applications, \file{run.bat} or the \file{raftcorr} command.
The main window opens empty: the image list shows the drop zone \ui{Drop image folder or Browse},
the canvas is blank and the progress area reads \ui{Ready}. The window does not shrink below
1120 $\times$ 760 pixels. The first time, the interface follows the language of the operating
system when RAFTcorr has a translation for it; change it in \menu{Settings}{Language}
(Section~\ref{sec:settings}).

\subsection{The models}
\label{sec:installing-models}

You need not install the models by hand. The first time an analysis uses a model that is not on
the computer, RAFTcorr downloads it from the release \file{weights-v1} of the RAFTcorr repository.
The \ui{PROGRESS} area says so, for example \emph{``Downloading model RAFTcorr-s4: 5.2 of 20.7
MB''}, and the bar follows the download; then the bar starts again from 0 for the tracking.
\ui{Cancel} stops the download and leaves no partial file. Every file is checked against its
SHA-256 before it is used, and kept for later runs; the \ui{LOG} names the address and the file
it is saved to. \ui{Suggest a Model} downloads s8 the same way when it is missing, and shows the download in
its progress dialog.

\screenshot[0.75\textwidth]{fig-models-window.png}{The Models window
(\menu{Help}{Models…}). One row per model with its downsampling factor, its reach, its state and
a \ui{Download} button while it is missing; below, the model folder, the link to the release page
and the buttons \ui{Download All}, \ui{Import Model Files…}, \ui{Open Model Folder} and
\ui{Close}.}{fig:models-window}

\menu{Help}{Models…} opens the Models window (Figure~\ref{fig:models-window}):

\begin{description}[leftmargin=0.6cm,style=nextline]
  \item[The model rows]
    Each row shows the model, \emph{``s = 4, reach 128 px''}, and its state: \ui{Installed} (the
    tooltip gives the path) or \ui{Not installed}, with the size of the download. A missing model
    has a \ui{Download} button.
  \item[Download All]
    Downloads every missing model, one after the other, and stops at the first failure.
    \ui{Cancel Download} appears while a download runs. Closing the window also cancels it.
  \item[Import Model Files…]
    Installs \file{.pth} files that you downloaded by hand. The SHA-256 of a file says which model
    it is, so a renamed download such as \file{RAFTcorr-s4 (1).pth} is recognised and stored under
    its release name; any other file is refused with \emph{``Not imported: \ldots\ Only the model
    files of the release can be imported.''}
  \item[Open Model Folder]
    Opens the per-user model folder, shown above the buttons as \ui{Model folder}.
\end{description}

When a run, or \ui{Suggest a Model}, needs a model that is neither on the computer nor
downloadable, the Models window opens instead of an error box. A card at the top names the model
and the reason, for example \emph{``RAFTcorr-s4 is needed, but it is not installed and could not
be downloaded.''} followed by \emph{``no connection to the download server (\ldots); check the
internet connection''}. Once the model is there, the card reads \emph{``RAFTcorr-s4 is installed
now. Close this window and start again.''} Then click \ui{Run DIC Analysis} again.

\subsubsection*{A computer without internet}

\begin{enumerate}
  \item On a computer with internet, open the release page of the models,
    \url{https://github.com/zachtong/RAFTcorr/releases/tag/weights-v1}; the Models window links
    to it. Download the files you need, from \file{RAFTcorr-s1.pth} to \file{RAFTcorr-s8.pth},
    about 20\,MB each.
  \item Copy them to the offline computer.
  \item Click \ui{Import Model Files…} in the Models window, or run
    \file{raftcorr models -{}-import FILE ...} (Section~\ref{sec:cli}).
\end{enumerate}

The Windows Setup program itself downloads PyTorch, so it cannot complete without internet. The
macOS disk image needs the connection only for the model files.

\subsubsection*{Where RAFTcorr looks for the model files}

For each model, RAFTcorr takes the first verified copy it finds, in this order:

\begin{enumerate}
  \item a file path that you give (\file{-{}-model} on the command line, or the Python library);
  \item the folder named by the environment variable \file{RAFTCORR\_WEIGHTS};
  \item the folder \file{models} of a source checkout;
  \item the per-user folder, \file{\%LOCALAPPDATA\%\textbackslash RAFTcorr\textbackslash models}
    on Windows, \file{\textasciitilde/.raftcorr/models} on macOS and Linux;
  \item a download into the per-user folder.
\end{enumerate}

A file whose SHA-256 does not match the release is skipped, never used, and the log says so:
\emph{``\ldots\ exists but its SHA-256 does not match the RAFTcorr-s4 release; skipping it''}.
\file{raftcorr models} lists the four models, where each one is installed and the folders
searched.

\subsection{Checking an installation}
\label{sec:installing-check}

\menu{Help}{Self-test} runs every installed model on a synthetic pair of $256 \times 256$ px
speckle images with a known translation, $u = +2.30$\,px and $v = -1.20$\,px, on the device set
in \ui{PARAMETERS}. The result opens in a dialog and goes to the log: one line per model with the
measured $u$ and $v$ and \file{ok} or \file{WRONG}, then \emph{``All installed models pass.''} or
\emph{``Some models FAILED.''} The self-test downloads nothing: a model that is not installed is
listed as \file{not installed}, and with no model installed at all the test fails.

From a terminal:

\begin{lstlisting}
raftcorr selftest              # the same test, on the device "auto"
raftcorr selftest --gui        # also open and close the main window offscreen
raftcorr info                  # versions, CUDA and the GPU
\end{lstlisting}

\file{raftcorr selftest} ends with exit code 0 when every check passes. \file{raftcorr info}
prints \file{CUDA available: True} or \file{False}, and with a GPU its name and free memory.
Section~\ref{sec:cli-other} describes both commands.

\subsection{GPU or CPU}
\label{sec:installing-gpu}

RAFTcorr runs on an NVIDIA GPU through CUDA, or on the CPU. For the GPU you need an NVIDIA card,
a current NVIDIA driver and a CUDA build of PyTorch. The Windows Setup program and
\file{install.bat} choose that build for you; with pip you install it yourself
(Section~\ref{sec:installing-pip}). On a Mac, RAFTcorr runs on the CPU. The \ui{Device} setting,
\ui{Automatic} by default, picks the GPU when PyTorch sees one (Section~\ref{sec:parameters});
\file{raftcorr info} tells you whether it does.
